\movedsub{What to do with a verbose generator}
\label{sec:deploy}

The detection rates above are low enough to raise a practical question the recommendations so far do not answer: given a generator that writes several assertions per answer, what should a system actually run? Our measurements support four steps, in order of how much they buy.

\textbf{Put a deterministic check in front, for anything symbolic.} That it flags every falsification we constructed is a consequence of how the probe defines them, not a measurement of the check. The measurement is what it costs --- 0.8 to 5.2\% of untouched answers flagged --- and what the learned scorers do on the same cases, which is 5.1\% to 89.4\%. A check indifferent to answer length is worth having where the aggregates are not, and it cannot see unsupported prose, so it is a pre-filter and not a gate.

\textbf{Do not put a fixed threshold on a long answer.} The scale floor moves with verbosity. Under the entailment aggregate, an untouched answer scored against an \emph{unrelated} document averages 0.245 for our terse arm on HotpotQA and 0.509 for the verbose one (Section~\ref{sec:verbosity}), so a 0.5 gate sits below the unrelated-document baseline for the verbose arm. Anchoring the gate to that floor costs one extra scoring pass over random contexts and makes the number mean the same thing across generators.

\textbf{If you control the generator, constrain the answer.} This is the largest single lever we measured, and it is free. One added prompt line moved Qwen from 1.77 to 4.90 assertions and detection from 35.1\% to 10.4\%; the same lever is available in the other direction. Shortening answers buys more detection than switching evaluators does.

\textbf{Choose the evaluator against the failure you care about.} Table~\ref{tab:choose} gives the properties. For separating two systems at a few hundred queries the binding constraint is resolving power, where AlignScore needs the smallest effect. For scoring a single verbose answer it is length robustness, where claim-min gives up the least --- with the caveat that its absolute scores fall near zero at five assertions, so it resists the bias without remaining thresholdable.

None of this is a general recipe. Every rate above is a property of a particular generator, evaluator and dataset, and the point of the paper is that those properties are not transferable. What does transfer is the procedure: falsify a value in your own outputs, score the floor against random contexts, and read the detection rate you actually have rather than the one the evaluator's benchmark reports.
